\subsection{Approaching is not the same as arriving}

\begin{corebox}[Two different questions]
The value $f(a)$ asks what happens at the input $a$. The limit
\[
\lim_{x\to a}f(x)
\]
asks what values are approached when $x$ is near $a$. These questions may have different answers.
\end{corebox}

\begin{examplebox}[A removable hole]
For
\[
f(x)=\frac{x^2-1}{x-1},\qquad x\ne1,
\]
we have $f(x)=x+1$ whenever $x\ne1$. Hence
\[
\lim_{x\to1}f(x)=2,
\]
even though $f(1)$ is not defined.
\end{examplebox}

\begin{definitionbox}[Function limit]
We write
\[
\lim_{x\to a}f(x)=L
\]
when the values $f(x)$ can be made arbitrarily close to $L$ by taking $x$ sufficiently close to $a$, without requiring $x=a$.
\end{definitionbox}

\begin{interpretationbox}[A limit describes behavior]
A limit is not primarily a substitution rule. It describes a stable pattern of nearby values. Substitution is only one possible method for recognizing that pattern.
\end{interpretationbox}
